\documentclass[11pt,a4paper]{article}
\usepackage[utf8]{inputenc}
\usepackage[T1]{fontenc}
\usepackage[french]{babel}
\usepackage{lmodern}
\usepackage[margin=2.2cm]{geometry}
\usepackage{booktabs,longtable,array,tabularx}
\usepackage[table]{xcolor}
\usepackage{enumitem}
\usepackage{hyperref}
\usepackage{fancyhdr}

\definecolor{ok}{HTML}{2E7D32}
\definecolor{part}{HTML}{E08E00}
\definecolor{ko}{HTML}{C62828}
\newcommand{\OK}{\textcolor{ok}{\textbf{Conforme}}}
\newcommand{\PART}{\textcolor{part}{\textbf{Partiel}}}
\newcommand{\KO}{\textcolor{ko}{\textbf{Absent}}}
\newcommand{\code}[1]{\texttt{\small #1}}

\pagestyle{fancy}\fancyhf{}
\lhead{SimHub -- Rapport à valider}\rhead{Septembre 2026}\cfoot{\thepage}
\setlist{nosep,leftmargin=1.5em}

\title{\textbf{SimHub}\\[0.3em]\large Gestion des droits par contexte Moodle et activité de cours\\ Audit de conformité au cahier des charges V0}
\author{Version 2 -- décisions du 26 septembre intégrées}
\date{26 septembre 2026}

\sloppy
\begin{document}
\maketitle

\section*{Synthèse}
\begin{itemize}
  \item \textbf{Constat.} Le plugin \code{local\_simhub} vérifie toutes ses capacités au
  \emph{contexte système} (51 appels à \code{context\_system::instance()}). Un enseignant
  inscrit dans son cours n'a donc aucun droit SimHub : un administrateur Moodle doit lui
  attribuer à la main un rôle système (\code{classes/local/roles.php}). Cela contredit le
  §11 (enseignant et responsable d'UC agissent sur \emph{leurs} enseignements) et le §10
  (« Cours / UC », « Permissions : gérer finement »).
  \item \textbf{Proposition (option 2 retenue).} Deux niveaux de contexte existants de
  Moodle, sans créer de niveau de contexte propre :
  \begin{enumerate}
    \item une \textbf{catégorie de cours SimHub} (une seule : chaque école a son propre
    Moodle) porte le référentiel transversal (ateliers, salle, ASV) : les rôles gestionnaire
    de salle et administrateur fonctionnel y sont attribués localement, par délégation, sans
    compte administrateur ;
    \item une \textbf{activité Moodle \code{mod\_simhub}} que le responsable d'UC ajoute
    lui-même dans son cours : les droits enseignant découlent automatiquement de
    l'inscription au cours, comme pour un devoir ou un test. Un enseignant n'a de pouvoir
    que sur les UC où il est inscrit comme enseignant, et en gagne sur une autre UC dès
    qu'il y est ajouté.
  \end{enumerate}
  \item \textbf{L'étudiant ne choisit jamais d'UC.} Il scanne le QR code de l'atelier et le
  réalise ; toutes les UC et tous les parcours qui contiennent cet atelier se mettent à jour
  seuls. C'est déjà le comportement du code actuel, et l'architecture le conserve.
  \item \textbf{Conformité.} Le socle V1 est largement couvert. Écarts principaux : droits
  enseignant (ci-dessus), carnet de notes et groupes Moodle non utilisés, filtres UC /
  parcours / statut personnel non exposés à l'étudiant, champ « catégorie » absent,
  export XLSX absent, contrôle de durée anti-faux scan absent.
\end{itemize}

\section{Problème : des capacités vérifiées au niveau système}

Dans Moodle, une permission se propage \emph{vers le bas} (système $\rightarrow$
catégorie $\rightarrow$ cours $\rightarrow$ activité), jamais vers le haut. La ligne
\code{'editingteacher' => CAP\_ALLOW} de \code{db/access.php} ne fixe que la valeur par
défaut du rôle ; elle ne sert que si ce rôle est attribué au contexte vérifié ou au-dessus.
Or SimHub vérifie au contexte système, où un enseignant n'a jamais de rôle.

Le correctif actuel (quatre rôles « Encadrant / Responsable d'UC / Gestionnaire de salle /
Administrateur fonctionnel SimHub », attribuables uniquement au niveau système) a trois
défauts :
\begin{itemize}
  \item chaque attribution passe par \emph{Administration du site $\rightarrow$ Attribuer des
  rôles système}, donc par un administrateur Moodle (\code{moodle/role:assign} au système) ;
  \item les droits sont \emph{globaux} : un responsable d'UC peut modifier les parcours de
  toutes les UC et suivre tous les étudiants, pas seulement les siens ;
  \item la liste des enseignants est tenue deux fois (inscriptions au cours et rôles
  système), ce que le §3.1 cherche justement à éviter.
\end{itemize}

\section{Architecture proposée}

\subsection{Répartition des responsabilités}

\begin{center}
\begin{tabularx}{\textwidth}{>{\raggedright}p{3.2cm}>{\raggedright}p{4.2cm}X}
\toprule
\textbf{Composant} & \textbf{Contexte des droits} & \textbf{Périmètre} \\
\midrule
\code{local\_simhub} (existant) & Catégorie de cours SimHub (réglage
\code{categoryid} unique) & Fiches ateliers, statuts, ressources, QR, plan,
import/export, grilles d'auto-évaluation, référentiel et pilotage ASV, accueil étudiant,
séances. \\
\addlinespace
\code{mod\_simhub} (nouveau) & Activité dans le cours de l'UC (hérite du cours) &
Rattachement d'ateliers à l'UC, parcours de l'UC, échéances, suivi des inscrits, validation
des séances de ses étudiants, note et achèvement dans le cours. \\
\addlinespace
Page publique ASV & Aucun (jeton signé + expiration) & Validation sur animal vivant par un
validateur externe (§9.3), inchangée. \\
\bottomrule
\end{tabularx}
\end{center}

Le lien entre les deux existe déjà en base : \code{local\_simhub\_rattachement.courseid},
\code{local\_simhub\_parcours.courseid} et \code{local\_simhub\_session.courseid}. L'activité
ne crée pas de nouveau référentiel : elle \emph{filtre} les données existantes sur son
cours. Un atelier reste unique et partagé entre toutes les UC (§6 « objet central »).

\subsection{Une séance d'atelier, plusieurs UC mises à jour}

La séance appartient à l'étudiant et à l'atelier, pas à une UC. L'étudiant scanne le QR
code, réalise l'atelier et s'auto-évalue, sans avoir à passer par le cours de l'UC. L'état
d'un atelier pour un étudiant est calculé à partir de ses seules séances sur cet atelier
(\code{parcours\_helper.php}, ligne 39), quels que soient l'UC ou le parcours enregistrés
sur la séance. Chaque activité \code{mod\_simhub} relit donc cet état pour les ateliers de
son UC et recalcule l'avancement et la note de ses inscrits.

\begin{center}
\begin{tabular}{l}
QR code de l'atelier $\rightarrow$ séance (étudiant, atelier) $\rightarrow$
\begin{tabular}[t]{@{}l@{}}UC 1 : avancement, note\\ UC 2 : avancement, note\\ parcours ASV,
etc.\end{tabular}
\end{tabular}
\end{center}

\noindent Les champs \code{session.courseid} et \code{session.parcoursid} restent des
indications de contexte facultatives, sans effet sur le calcul. À préserver :
\code{grade\_update()} doit être déclenché pour \emph{chaque} activité dont l'UC contient
l'atelier, lorsqu'une séance se termine ou est validée (observateur de l'événement
\code{session\_completed}).

\subsection{Profils du §11 et attribution des droits}

\begin{center}
\begin{tabularx}{\textwidth}{>{\raggedright}p{3.3cm}>{\raggedright}p{4.5cm}X}
\toprule
\textbf{Profil} & \textbf{Rôle Moodle et contexte} & \textbf{Qui l'attribue} \\
\midrule
Étudiant & Utilisateur authentifié (accueil) + étudiant du cours (activité) &
Automatique (connexion, inscription au cours). \\
Enseignant / formateur & \code{teacher} dans le cours de l'UC & Inscription habituelle au
cours, aucun geste SimHub. \\
Responsable d'UC & \code{editingteacher} dans le cours de l'UC & Idem ; il ajoute
lui-même l'activité (\code{moodle/course:manageactivities}). \\
Gestionnaire de salle & Rôle SimHub attribué \emph{dans la catégorie} & L'administrateur
fonctionnel. \\
Administrateur fonctionnel & Rôle SimHub attribué dans la catégorie, avec
\code{moodle/role:assign} limité à celle-ci & Un administrateur Moodle, une seule fois. \\
Validateur externe ASV & Aucun compte & Lien ou QR code à jeton. \\
\bottomrule
\end{tabularx}
\end{center}

\noindent\textbf{Réponse à la question initiale :} seule la désignation de
l'administrateur fonctionnel passe encore par un compte administrateur
Moodle. Tout le reste est délégué : les enseignants par les inscriptions aux cours, les
gestionnaires de salle par l'administrateur fonctionnel dans sa catégorie.

\subsection{Capacités}

\begin{center}
\begin{tabularx}{\textwidth}{>{\raggedright}p{6.4cm}>{\raggedright}p{2.6cm}X}
\toprule
\textbf{Capacité} & \textbf{Contexte} & \textbf{Archétypes par défaut} \\
\midrule
\code{local/simhub:view}, \code{startsession}, \code{submitautoeval} & Catégorie & user \\
\code{local/simhub:manageateliers}, \code{manageressources}, \code{managestatuts},
\code{manageqrcodes}, \code{importexport} & Catégorie & manager \\
\code{local/simhub:manageasv}, \code{configure} & Catégorie & manager \\
\code{local/simhub:validateasvsimulation} & Catégorie \emph{et} cours & editingteacher,
teacher \\
\addlinespace
\code{mod/simhub:addinstance} & Cours & editingteacher, manager \\
\code{mod/simhub:view} & Activité & student, teacher, editingteacher \\
\code{mod/simhub:viewprogression}, \code{validatesession}, \code{exportsuivi} & Activité &
teacher, editingteacher \\
\code{mod/simhub:manageparcours}, \code{managerattachement} & Activité & editingteacher \\
\bottomrule
\end{tabularx}
\end{center}

Les capacités \code{local/simhub:viewprogression}, \code{validatesession},
\code{exportsuivi}, \code{manageparcours} et \code{managerattachement} restent au niveau
catégorie pour les gestionnaires et l'administrateur fonctionnel (vue transversale).

\subsection{Ce que l'activité apporte en plus}
\begin{itemize}
  \item \textbf{Carnet de notes} (§10, §3.1) : \code{grade\_update()} publie l'avancement du
  parcours de l'UC en pourcentage d'avancement (0 à 100) dans le carnet du cours.
  \item \textbf{Achèvement d'activité} : réutilisable par l'achèvement de cours et les badges
  de cours Moodle, sans code de badge spécifique.
  \item \textbf{Groupes} (§10) : le mode de groupe standard de l'activité permet à un
  encadrant de ne suivre que son groupe de TP.
  \item \textbf{Sauvegarde / restauration} et réinitialisation annuelle du cours :
  l'activité suit le cours d'une promotion à l'autre.
\end{itemize}

\subsection{Migration}
\begin{enumerate}
  \item Réglage \code{categoryid} unique ; un utilitaire
  \code{local\_simhub\textbackslash local\textbackslash contexte::racine()} remplace les 51
  appels à \code{context\_system::instance()}.
  \item Les rôles SimHub deviennent attribuables au niveau catégorie
  (\code{set\_role\_contextlevels}). \textbf{Compatibilité :} une attribution système
  existante reste valable, puisque le système est parent de la catégorie.
  \item Les rôles « Encadrant SimHub » et « Responsable d'UC SimHub » deviennent inutiles
  pour l'usage courant ; ils sont conservés pour une vue transversale.
  \item Une étape de \code{db/upgrade.php} ajoute \code{local\_simhub\_parcours.cmid}
  (nullable) pour lier un parcours à l'instance d'activité qui l'a créé.
  \item Nouveau plugin \code{mod/simhub} : \code{mod\_form.php}, \code{lib.php}
  (\code{simhub\_add\_instance}, \code{supports}, notes), \code{view.php}, sauvegarde.
\end{enumerate}

\subsection{Décisions validées}
\begin{center}
\begin{tabularx}{\textwidth}{>{\raggedright}p{5.2cm}X}
\toprule
\textbf{Question} & \textbf{Décision et conséquence} \\
\midrule
Découpage par établissement & Aucun : chaque école a son propre Moodle. Une seule
catégorie SimHub. \code{envcode} ne sert plus qu'à la \textbf{traçabilité} (origine des
données dans les exports, imports et PDF), jamais au filtrage ni aux droits. \\
\addlinespace
Portée des droits d'un enseignant & Uniquement les UC où il est inscrit comme enseignant.
Il valide les séances d'un étudiant inscrit à son UC portant sur un atelier rattaché à son
UC, où que la séance ait été lancée. Ajouté à une autre UC, il y obtient les mêmes droits. \\
\addlinespace
Grille d'auto-évaluation par le responsable d'UC & Oui, pour les ateliers rattachés à son
UC. La grille étant partagée par toutes les UC, chaque modification est tracée
(auteur, date) et visible sur la fiche atelier. \\
\addlinespace
Note au carnet & Pourcentage d'avancement du parcours de l'UC. \\
\addlinespace
Validation ASV en simulation & Les deux : les enseignants dans leur cours \emph{et} des
formateurs désignés au niveau de la catégorie SimHub. \\
\bottomrule
\end{tabularx}
\end{center}

\section{Audit de conformité au cahier des charges}

Vérification faite dans le code (schéma \code{db/install.xml}, pages, classes), et non à
partir du seul README. Légende : \OK, \PART, \KO.

\renewcommand{\arraystretch}{1.15}
\begin{longtable}{>{\raggedright}p{1.3cm}>{\raggedright}p{4.2cm}p{1.95cm}>{\raggedright\arraybackslash}p{7.05cm}}
\toprule
\textbf{§} & \textbf{Exigence} & \textbf{État} & \textbf{Constat} \\
\midrule\endhead
\bottomrule\endfoot
4 & Moodle comme socle (rôles, cohortes, groupes, UC) & \PART & Cohortes et cours utilisés ;
rôles détournés au niveau système ; groupes jamais utilisés. \\
4 & Paramétrable ENVF & \OK & Chaque école a son Moodle : référentiels, logo et nom
propres. \code{envcode} conservé pour la seule traçabilité. \\
7 & Scan sans passer par l'UC & \OK & L'état d'un atelier ne dépend que des séances de
l'étudiant sur cet atelier ; toutes les UC et tous les parcours en profitent. \\
4 & Mobile d'abord & \PART & Cartes et grille verticale prévues ; pas de test sur écran de
téléphone documenté. \\
5.1 & Accueil personnalisé (7 sections) & \OK & Les sept sections sont présentes
(\code{student\_home\_page.php}). \\
5.2 & Filtres UC, parcours & \PART & Gérés côté serveur (\code{atelier\_filter.php}, lignes
120--126) mais \emph{absents du formulaire} de l'accueil. \\
5.2 & Filtre statut personnel & \KO & Aucun filtre « pas commencé / réalisé / validé / à
reprendre » ; l'état n'est qu'affiché sur la carte. \\
5.2 & Filtre statut atelier & \PART & Paramètre \code{statut} accepté, non exposé ; le défaut
(actifs + indisponibles) est conforme. \\
5.2 & Discipline, espèce, niveau, durée, année, mot-clé & \OK & Présents dans le formulaire. \\
5.3 & Carte : « Catégorie » & \KO & Aucun champ catégorie dans
\code{local\_simhub\_atelier}. \\
5.3 & Carte : UC associées et niveau attendu & \PART & \code{rattachement.niveauattendu}
existe ; affichage sur la carte à vérifier. \\
5.3 & Numéro invariant, boutons localisation / ressources / commencer & \OK & \\
5.4 & Localisation, plan avec repère cliquable & \OK & \code{atelier\_plan.php}. \\
5.5, 6.2 & Ressources étudiantes / sources éditables & \OK & Source éditable jamais servie
à l'étudiant. \\
5.6, 7.2 & Auto-évaluation 3 niveaux, rubrique risques, auto-bilan & \OK & \\
6.2 & Grille paramétrable par les enseignants & \PART & Réservée à
\code{manageateliers} (gestionnaire de salle). Voir point à valider 3. \\
6 & Fiche : identification, localisation, statut, référent, commentaire & \OK & Nom
court/long, discipline, espèce, niveau, durée, salle/zone/poste, référent. \\
6 & Rattachement : UC, année, cohorte, recommandé/obligatoire & \OK & Table
\code{rattachement} complète. \\
6.1 & Statuts et indisponibilité (motif, référent, date, échéance) & \OK & Historique
\code{indispo}. \\
7 & QR code relié à la fiche, démarrage de séance & \OK & Génération locale, sans service
tiers. \\
7.1 & Suivi : étudiant, atelier, date, UC/parcours, statut & \OK & Table \code{session}. \\
7.1 & Temps passé pour détecter un contournement & \PART & \code{timestart}/\code{timeend}
enregistrés, mais aucune durée minimale ni signalement des séances anormalement courtes. \\
7.3 & Code de séance, validation encadrant non bloquante & \OK & \\
7.3 & Reconnaissance réseau local & \KO & Non implémentée (V1+, dépend de l'infrastructure
Wi-Fi). \\
8 & Types de parcours, ordre, obligatoire, échéances & \OK & \\
8.1 & Suivi enseignant : non commencés, en retard, à reprendre, par cohorte / UC & \PART &
Présent (\code{dashboard.php}) mais non restreint aux étudiants de l'enseignant ; pas de
vue par groupe. \\
9 & Module ASV complet (simulation, animal vivant, livret, A1--A3) & \OK & Ordre imposé,
signature obligatoire, annulation tracée, certification globale A3. \\
10 & Rôles : différencier les profils & \PART & Profils distincts, mais au niveau système
(section 1). \\
10 & Cours / UC : rattacher aux enseignements & \PART & Rattachement par \code{courseid},
mais aucune présence de SimHub dans le cours lui-même. \\
10 & Cohortes et groupes & \PART & Cohortes oui, groupes non. \\
10 & Carnet de notes & \KO & Aucun appel à l'API des notes. \\
10 & Badges / attestations & \OK & Badges de site existants, attestations PDF. \\
11 & Enseignant : \emph{ses} enseignements & \KO & Droits globaux, attribués par un
administrateur. \\
12.1 & Import souple (ateliers, rattachements, parcours) & \PART & CSV avec alias de
colonnes ; pas d'import XLSX ; import des salles/zones et des ressources non couvert en tant
que tel. \\
12.2 & Tableaux de bord étudiant, enseignant, salle, ASV & \OK & \\
12.3 & Exports CSV, XLSX, PDF & \PART & CSV et PDF ; \textbf{XLSX absent} alors que
l'API \code{dataformat} de Moodle le fournit. \\
12.3 & Exports par étudiant, UC, parcours, cohorte & \PART & Par parcours et liste des
ateliers ; pas d'export direct par UC ni par cohorte. \\
14 & Hors périmètre respecté & \OK & Pas de ticketing, OSCE ni signature qualifiée. \\
\end{longtable}

\noindent\textbf{Remarque mineure :} \code{version.php} indique encore « non testé sur
instance réelle », alors que le README documente des essais sur Moodle~4.5 et~5.0.

\section{Plan de travail proposé}

\begin{center}
\begin{tabularx}{\textwidth}{>{\raggedright}p{0.6cm}>{\raggedright}p{7.2cm}>{\raggedright}p{2cm}X}
\toprule
\textbf{N°} & \textbf{Lot} & \textbf{Phase} & \textbf{Dépend de} \\
\midrule
1 & Contexte catégorie : réglage, utilitaire de contexte, rôles attribuables en catégorie ;
\code{envcode} retiré des filtres & V1 & -- \\
2 & Activité \code{mod\_simhub} : ajout au cours, parcours et rattachements de l'UC, suivi des
inscrits & V1 & Lot 1 \\
3 & Carnet de notes (\% d'avancement) et achèvement, mis à jour à chaque séance & V1 & Lot 2 \\
4 & Filtres UC, parcours et statut personnel sur l'accueil ; champ catégorie & V1 & -- \\
5 & Signalement des séances trop courtes (durée minimale paramétrable) & V1+ & -- \\
6 & Export XLSX via \code{dataformat} ; exports par UC et par cohorte & V1 & -- \\
7 & Groupes dans le suivi enseignant & V1 & Lot 2 \\
8 & Reconnaissance réseau local & V1+ & Infrastructure des ENV \\
\bottomrule
\end{tabularx}
\end{center}

\end{document}
